\subsection{Trees, Forests, and Boosting}

Tree-based methods partition the input space into simple regions and then combine many weak or unstable predictors into a stronger rule.

\begin{definition} \label{def:sl-decision-tree} \hlt{Decision Tree}\\
A decision tree recursively splits the input space using rules such as $x_j\leq t$.  Each leaf stores a prediction, commonly a class distribution or an average response.
\end{definition}

\begin{definition} \label{def:sl-bagging-random-forest} \hlt{Bagging and Random Forests}\\
\hlt{Bootstrap aggregating} trains predictors on bootstrap resamples and averages their outputs.  A \hlt{random forest} adds random feature selection at tree splits, reducing correlation among the trees.
\end{definition}

\begin{definition} \label{def:sl-boosting} \hlt{Boosting}\\
Boosting builds an additive predictor
\[
F_M(x)=\sum_{m=1}^M\alpha_m h_m(x)
\]
by fitting each weak learner $h_m$ to the errors, residuals, or weighted examples emphasized by the current ensemble.
\end{definition}

\begin{remark} \label{rem:sl-ensemble-source-migration} \hlt{Ensemble Methods}\\
The old classical-machine-learning source contained only an ensemble heading.  The substantive bagging and boosting notes in the deep-learning regularization source have been migrated here because they are general ensemble methods, not neural-network-specific constructions.
\end{remark}
